\documentclass[11pt, oneside]{article}   	% use "amsart" instead of "article" for AMSLaTeX format
\usepackage{geometry}                		% See geometry.pdf to learn the layout options. There are lots.
\geometry{letterpaper}                   		% ... or a4paper or a5paper or ... 
%\geometry{landscape}                		% Activate for rotated page geometry
%\usepackage[parfill]{parskip}    		% Activate to begin paragraphs with an empty line rather than an indent
\usepackage{graphicx}				% Use pdf, png, jpg, or eps§ with pdflatex; use eps in DVI mode
								% TeX will automatically convert eps --> pdf in pdflatex		
\usepackage{amssymb}

%SetFonts

%SetFonts


\title{Useful words and turns of phrase for paper}
\author{Jacqueline McCleary}
%\date{}							% Activate to display a given date or no date

\begin{document}
\maketitle
%\section{}
%\subsection{}

From Menard 2010: 

``Our observational results primarily make use of the angular cross-correlation between the brightness of quasars and the projected density of galaxies'' 

``Secondly, this large sample – together with high accuracy photometry in five optical pass bands – allows us to detect the presence of dust in the intervening space and explore its distribution and properties. This allows us to study the properties of intergalactic dust, and provides a way of inferring the abundance of dust in the Universe.''

``On small scales ($<0.01$), the signal contributions from magnification and dust extinction are expected to be dominated by Poisson noise. For larger scales, measurements from different angular bins become significantly correlated, as one would expect since different quasars sample the same local population of galaxies.'' \textbf{What to do about this isn't clear}. I *still* think this is what's behind the fat tail, though I feel like I tested this before and never got a satisfactory answer. Oh right, the Y1 mask was effectively a test of this. It might be worth trying again... or I could just leave it alone. 

From ARA\&A Cosmic Metals paper:

``Although their mass grows with time, only a minority of the baryonic matter is found in stars; even today, some 13.7 Gyr after the Big Bang, $>90\%$ 
of the baryons are found in the gaseous phase of the Universe. The gas, notably the cold gas traced by HI and H$_2$, provides the reservoir of fuel for forming stars.''

\end{document}  